% !tp-example Basic usage
\documentclass{article}
\begin{document}
\wcTpGraphGeometry{}
\end{document}

% !tp-example Javascript API
\documentclass{article}
\begin{document}
\wcScript[type=module]|
import "/components/graph-geometry/graph-geometry.js";

const example = document.createDocumentFragment();

const element1 = document.createElement("tp-graph-geometry");
example.append(element1);

document.body.append(example);
|
\end{document}

% !tp-example Attached points
\documentclass{article}
\begin{document}
\wcTpGraphGeometry{
  \wcScript[type=tp/geometry]|
point O = (0, 0)
point A = (6, 0)
circle c = circle(O, 5)
point P = on(c, angle=60deg)
line d = line(O, A)
point Q = on(d, t=0.5)
|
}
\end{document}

% !tp-example Arc only
\documentclass{article}
\begin{document}
\wcTpGraphGeometry{
  \wcScript[type=tp/geometry]|
point O = (0, 0)
circle c = circle(O, 5) hidden
arc a = arc(c, start=20deg, end=140deg)
|
}
\end{document}

% !tp-example From a file
\documentclass{article}
\begin{document}
\wcTpGraphGeometry[src=/docs/components/graph-geometry/examples/circle.geometry]{}
\end{document}

% !tp-example Ruler primitives
\documentclass{article}
\begin{document}
\wcTpGraphGeometry[grid]{
  \wcScript[type=tp/geometry]|
point A = (-4, 0)
point B = (0, 3)
point O = (4, 0)
segment s = segment(A, B)
segment fixed = segment(A, B, length=3)
point E = on(fixed, t=1)
line d = line(B, O)
ray r = ray(A, O)
vector v = vector(O, B)
circle c = circle(O, 2)
circle through = circle(A, through=B) hidden
sector w = sector(c, start=20deg, end=100deg)
comment note = comment(E, "AE = 3 units")
|
}
\end{document}

% !tp-example Direct constructions
\documentclass{article}
\begin{document}
\wcTpGraphGeometry[built-in,grid]{
  \wcScript[type=tp/geometry]|
point A = (-4, -2)
point B = (2, -2)
point O = (0, 1)
point P = (5, 4)
segment s = segment(A, B)
circle c = circle(O, 2)
point M = midpoint(s)
line m = perpendicularBisector(s)
line p = parallel(O, s)
line q = perpendicular(A, s)
line b = bisector(A, O, B)
line t0 = tangent(P, c, branch=0)
line t1 = tangent(P, c, branch=1)
comment note = comment(P, "Move P")
|
}
\end{document}

% !tp-example Styles and export
\documentclass{article}
\begin{document}
\wcTpGraphGeometry[grid]{
  \wcScript[type=tp/geometry]|
point O = (0, 0)
point A = (-3, -2)
point B = (3, 2)
segment s = segment(A, B) style=dashed color=blue
circle c = circle(O, 3) style=dotted color=red
vector v = vector(O, B) color=green
circle hiddenCircle = circle(A, 2) style=hidden color=violet
comment note = comment(A, "Select an object to style it") color=blue
|
}
\end{document}

% !tp-example Polygons
\documentclass{article}
\begin{document}
\wcTpGraphGeometry[built-in,grid]{
  \wcScript[type=tp/geometry]|
point A = (-5, -2)
point B = (-2, -2)
point C = (-3, 0)
point D = (-5, 2)
point E = (1, -2)
point F = (4, -2)
polygon irregular = polygon(A, B, C, D) color=blue
polygon regular = regularPolygon(E, F, sides=5) color=green
|
}
\end{document}

% !tp-example Transformations
\documentclass{article}
\begin{document}
\wcTpGraphGeometry[built-in,grid]{
  \wcScript[type=tp/geometry]|
point O = (0, 0)
point X = (0, 3)
point A = (1, -1)
point B = (2, -1)
point C = (1, 0)
point T = (3, 2)
line axis = line(O, X) style=dashed color=zinc
vector v = vector(O, T) color=zinc
polygon original = polygon(A, B, C) color=orange
polygon reflected = axialSymmetry(original, axis) color=blue
polygon opposite = centralSymmetry(original, O) color=violet
polygon enlarged = homothety(original, O, ratio=2) color=green
polygon turned = rotation(original, O, angle=90deg) color=red
polygon shifted = translation(original, v) color=teal
|
}
\end{document}

% !tp-example Transform in place
\documentclass{article}
\begin{document}
\wcTpGraphGeometry[built-in,grid]{
  \wcScript[type=tp/geometry]|
point O = (0, 0)
point A = (3, 0)
circle c = circle(A, 1) color=blue
point Q = on(c, angle=30deg)
|
}
\end{document}

% !tp-example Measurements
\documentclass{article}
\begin{document}
\wcTpGraphGeometry[grid]{
  \wcScript[type=tp/geometry]|
point A = (-4, -1)
point B = (0, -1)
point C = (0, 2)
segment AB = segment(A, B) color=blue
segment BC = segment(B, C) color=green
polygon triangle = polygon(A, B, C) color=blue
point O = (4, 1)
circle c = circle(O, 2) color=orange
comment side = length(AB) color=blue
comment corner = angle(AB, BC) color=green
segment AC = segment(A, C) color=blue
comment slope = angle(AB, AC) color=blue
comment surface = area(c) color=orange
|
}
\end{document}

% !tp-example Interactive figure
\documentclass{article}
\begin{document}
\wcTpGraphGeometry[label=Circle and attached points,canvas-only,grid,src=/docs/components/graph-geometry/examples/circle.geometry]{}
\end{document}

% !tp-example Shared points and vectors
\documentclass{article}
\begin{document}
Open the source and change A to (6, 2), then press Play. Vector u follows the two endpoints. The same source works in both plot components.


\wcTpGraphGeometry[canvas-only,style=--tp-plot-height:440px]{
  \wcScript[type=tp/geometry]|
label = "Points and vectors"
point O = (0, 0)
point A = (3, 4)
vector u = vector(O, A)
|
}
\end{document}
